\section{Introduction}\label{sec:introduction}

Let $K$ be an imaginary quadratic number field and let $E/K$ be an
elliptic curve. The automorphy conjecture for $E$ asks for an automorphic
representation of $\GL_2(\A_K)$ whose local parameters realize the
cohomology of $E$. At a place of good reduction this requires the Hecke
polynomial to recover the point count of the reduced curve; at a bad
place it requires the full local parameter, including its inertia and
monodromy. At infinity the parameter is prescribed by the Hodge types
$(1,0)$ and $(0,1)$.

We prove this conjecture for every $K$ and $E$. We use unitary
normalization on the automorphic side and the Tate-normalized local
Langlands correspondence on the cohomological side. These conventions
are recalled in Section~\ref{sec:conventions}.

\begin{theorem}\label{thm:main}
Let $K$ be an imaginary quadratic number field and $E/K$ an elliptic
curve. There is an isobaric automorphic representation $\pi_E$ of
$\GL_2(\A_K)$ with trivial central character and weight zero such that,
for every prime $l$ and every isomorphism
$\iota:\overline\Q_l\xrightarrow{\sim}\C$,
\[
 r_\iota(\pi_E)\simeq
 H^1_{\et}(E_{\overline K},\overline\Q_l).
\]
At every finite place $v\nmid l$ one has
\[
 \iota\WD\bigl(H^1_{\et}(E_{\overline K},\overline\Q_l)
                     |_{G_{K_v}}\bigr)^{\Fss}
 \simeq
 \rec_{K_v}\bigl(\pi_{E,v}\otimes|\det|_v^{-1/2}\bigr).
\]
The unitary parameter at the complex place is
$z\mapsto\diag(z/|z|,\overline z/|z|)$.
In particular every local factor agrees, and
\[
 L(E/K,s)=L(\pi_E,s-1/2).
\]
If $E$ has no complex multiplication, $\pi_E$ is cuspidal. In the CM
case it is given by automorphic induction of a Hecke character, or by
an isobaric sum when the CM field is already $K$.
\end{theorem}

Here $r_\iota$ in the CM isobaric case is the sum of the corresponding
character realizations with the same Tate normalization.

\subsection{Historical context}

The arithmetic groups associated with imaginary quadratic fields and
their action on hyperbolic three-space originate in Bianchi's
work~\cite{Bianchi}. Their presentations and cohomology were developed
further by Swan and Harder~\cite{Swan,Harder}. The expected relationship
with elliptic curves belongs to the Langlands program: in
\cite[Section~7]{LanglandsProblems}, Langlands explicitly discusses the
automorphy of local representations attached to elliptic curves over
global fields of characteristic zero. Grunewald--Helling--Mennicke,
Cremona, and Cremona--Whitley developed substantial computations relating
elliptic curves and automorphic forms over imaginary quadratic
fields~\cite{GHM,Cremona,CremonaWhitley}.

Over $\Q$, Wiles and Taylor--Wiles proved modularity of semistable
elliptic curves~\cite{Wiles,TaylorWiles}, and
Breuil--Conrad--Diamond--Taylor proved modularity of every elliptic
curve~\cite{BCDT}. A central method is to lift automorphy from a residual
Galois representation. Wiles's $3$--$5$ switch also shows how an auxiliary
elliptic curve, chosen with prescribed torsion and local behavior, can
replace a residual representation that does not satisfy the initial
lifting hypotheses.

For imaginary quadratic fields, the relevant cohomology occurs in more
than one degree. The construction of Galois representations and the
lifting problem both require additional work. Harris--Soudry--Taylor and
Taylor constructed representations attached to Bianchi forms in important
cases~\cite{HST,Taylor1994}. For conjugation-invariant central character,
Berger--Harcos obtained Frobenius compatibility away from finitely many
primes~\cite{BergerHarcos}. Harris--Lan--Taylor--Thorne constructed
Galois representations for general regular algebraic cuspidal
representations over CM fields~\cite{HLTT}, and Scholze treated torsion
cohomology as well~\cite{Scholze}. Calegari--Geraghty developed patching
for cohomology in several degrees~\cite{CalegariGeraghty}.
The potential automorphy theorem of
Allen--Calegari--Caraiani--Gee--Helm--Le Hung--Newton--Scholze--Taylor--Thorne
implies potential modularity of every elliptic curve over a CM
field~\cite[Theorem~1.0.1]{ACC}. Boxer--Calegari--Gee--Pilloni also proved
potential modularity for genus-one curves over quadratic extensions of
totally real fields~\cite[Theorem~1.1.4]{BCGP}. Potential modularity permits
an extension of the ground field; automorphy over the original field is a
further problem.

Allen--Khare--Thorne proved residual modularity results in small
characteristics and deduced modularity for a positive proportion of
elliptic curves over each CM field not containing a primitive fifth root
of unity~\cite{AKT}. Caraiani--Newton proved a lifting theorem allowing
small coefficient primes and ramified CM fields, together with the
local-global compatibility needed for its applications~\cite{CN}.
Their results give modularity under suitable residual-image hypotheses
in characteristic $3$ or $5$, modularity for $100\%$ of Weierstrass
equations, ordered by height, over each Galois CM field not containing $\zeta_5$, and
modularity of every elliptic curve over an imaginary quadratic field
$K$ for which $X_0(15)(K)$ has rank zero
\cite[Theorem~6.1 and Corollaries~6.1.2 and~7.1.2]{CN}.
Theorem~\ref{thm:main} removes the remaining restrictions on the curve
and the imaginary quadratic field. Caraiani--Newton's lifting theorem
and residual-image criterion for modularity of elliptic curves remain
essential inputs.

\subsection{The geometric prime switch}

The difficulty is to reach a residual representation to which the known
modularity theorem applies while retaining control at the primes where
automorphy will be lifted. We construct the intermediate representations
as character pieces of a single family of Jacobians.

Fix a non-CM elliptic curve $E/K$. Choose a sufficiently large prime $p$
for which the image on $E[p]$ contains $\SL_2(\F_p)$, and put
$n=3p$ and $F=K(\mu_n)$. For $t$ in a nonempty affine open over $K$, we
construct a genus-$n$ cyclic cover $C_n(t)$ of a quadratic twist $B_t$ of
$E$. It is branched at two points and admits an involution that inverts
the rotations. For a character $\theta$ of the rotation group $\mu_n$,
write
\[
 V_\theta=H^1(C_n(t))_\theta.
\]
Each $V_\theta$ is a rank-two compatible system with determinant
$\eps_l^{-1}$, where $\eps_l$ is the $l$-adic cyclotomic character, and
labeled Hodge--Tate weights $0,1$.

Choose characters $\eta$ and $\psi$ of orders $3$ and $p$, respectively,
and set $\chi=\eta\psi$. The order-three system $V_\eta$ comes from an
elliptic curve $A/F$, while $V_1=H^1(B_t)$. An integral description of the
cohomology gives the chain
\begin{equation}\label{intro:chain}
 \underbrace{V_\eta}_{H^1(A)}
 \ \equiv\!\!\pmod p\ \ V_\chi
 \ \equiv\!\!\pmod 3\ \ V_\psi
 \ \equiv\!\!\pmod p\ \ \underbrace{V_1}_{H^1(B_t)}.
\end{equation}
Each congruence means equality of the semisimplified residual Galois
representations. The order-$p$ part of a character disappears modulo
$p$, and the order-$3$ part disappears modulo $3$. Automorphy will pass
from left to right. The prime $5$ is used only to prove the initial
modularity of $A$.

There are two useful degenerations of the covers. At one, the Jacobian
has compact-type limit $E^n$. This yields a free integral group-ring
description of its cohomology and lets us impose the local reduction
behavior of all the character pieces at once. At the other, explicit
vanishing cycles give opposite unipotent elements on each nontrivial
rank-two piece. The latter supply absolute irreducibility and suitable
Frobenius elements in residual characteristic $3$, $5$, or $p$.
The two degeneration pictures are given in
Section~\ref{sec:geometry}, and the vanishing-cycle calculation in
Section~\ref{sec:monodromy}.

Two specialization arguments complete the construction. First, finite
torsion is locally constant in a family over a local field. Once
semistability is secured, the full Frobenius polynomials belong to a
finite set. Sufficiently deep torsion comparison with the fiber $E^n$
therefore preserves its toric rank and, in the good-reduction case, its
Newton slopes. This provides simultaneous local conditions for every
character piece. Second, Hilbert irreducibility is applied to the
restriction of scalars of the parameter line from $K$ to $\Q$. The two
conjugate families then occur in one finite monodromy quotient. This
ensures the required residual genericity at every place above a rational
split prime, while retaining the prescribed local conditions on $t$.

The construction has close antecedents. Cyclic covers of elliptic
curves with two branch points and dihedral symmetry occur in the
real-multiplication constructions of Tautz--Top--Verberkmoes and in
Ellenberg's treatment of Jacobian endomorphism
algebras~\cite{TTV,Ellenberg}; Mestre's constructions form part of the
same geometric tradition~\cite{Mestre}. Darmon developed congruences that
remove prime-power parts of monodromy and a conditional modularity
argument based on successive lifts for rigid local
systems~\cite{Darmon1999,Darmon2000}.
The coefficient $2-\zeta-\zeta^{-1}$ in our monodromy calculation also
appears in the rank-two monodromy studied by
Darmon--Mestre~\cite{DarmonMestre}. Taylor's use of auxiliary
abelian varieties with real multiplication is another higher-dimensional
form of prime switching~\cite{Taylor2002}.

The feature developed here is the combination for a \emph{prescribed}
elliptic curve: the character orders $1,p,3p,3$ occur in one integral
cohomology lattice, an elliptic curve occurs at each end of the chain,
and the compact-type degeneration and simultaneous residual specialization
make every step applicable to the CM-field lifting theorem. The local
constancy argument and the two-conjugate specialization argument are
formulated separately, so they can be used independently of the final
automorphy application.

Sections~\ref{sec:geometry} and~\ref{sec:monodromy} establish the cover
geometry and integral monodromy. Section~\ref{sec:systems} constructs the
rank-two systems and their congruences. Section~\ref{sec:inputs} states the
automorphy inputs at the point where their hypotheses become concrete
conditions on those systems. Sections~\ref{sec:local}
and~\ref{sec:residual} find one parameter satisfying those conditions.
Section~\ref{sec:automorphy} performs the three lifts, descends from $F$
to $K$, and completes the local comparisons and the CM case.
